% ============================================================
% DIAGRAM TEMPLATE – From Biological Neurons to LLMs
% ============================================================
% Usage:
%   1. Copy this template to book/tikz/{paper_id}_{diagram_name}.tex
%   2. Fill in the title, caption, and diagram code
%   3. Compile with: pdflatex {filename}.tex
%   4. Convert to SVG: pdf2svg {filename}.pdf {filename}.svg
% ============================================================

\documentclass[tikz,border=10pt]{standalone}
\usepackage{tikz}
\usepackage{amsmath}
\usepackage{amsfonts}

% ============================================================
% Colour Palette
% ============================================================
% Neuron body:       blue!20
% Input:             green!20
% Output:            red!20
% Weight:            orange!30
% Bias/Threshold:    gray!20
% Edge/Arrow:        black!70
% Text:              black
% Highlight:         yellow!20

\usetikzlibrary{shapes,arrows,positioning,fit,calc,backgrounds}

\begin{document}

% ============================================================
% DIAGRAM: [Title]
% ============================================================
% Caption: [Figure X.Y: Description. (Source: Paper, Year)]
% ============================================================

\begin{tikzpicture}[
    % Node styles
    neuron/.style={draw, circle, minimum size=1cm, inner sep=0pt, fill=blue!20},
    input/.style={draw, circle, minimum size=0.8cm, inner sep=0pt, fill=green!20},
    output/.style={draw, circle, minimum size=0.8cm, inner sep=0pt, fill=red!20},
    weight/.style={draw, rectangle, minimum width=0.6cm, minimum height=0.4cm, fill=orange!30, font=\tiny},
    threshold/.style={draw, rectangle, minimum width=0.6cm, minimum height=0.4cm, fill=gray!20, font=\tiny},
    label/.style={font=\sffamily\small},
    title/.style={font=\sffamily\Large\bfseries},
    caption/.style={font=\sffamily\small\itshape},
    edge/.style={-stealth, thick, draw=black!70},
]

% ============================================================
% Diagram Code Here
% ============================================================

\end{tikzpicture}

\end{document}